% ============================================================
\chapter{Risk Measures}
\label{ch:risk}
% ============================================================

\section{Introduction}

After the financial crises of the 1990s --- the Orange County debacle (1994), the Barings collapse (1995), the LTCM crisis (1998) --- the question of measuring risk became central. J.P.\,Morgan popularized \emph{Value at Risk} (VaR) in 1994 with its RiskMetrics system: a single number summarizing the maximum probable loss over a given horizon at a given confidence level. But Philippe Artzner and colleagues showed in 1999 that VaR is not \emph{coherent}: it violates subadditivity, meaning that diversifying a portfolio can, paradoxically, increase measured VaR. \emph{Expected Shortfall} (CVaR) corrects this flaw. This chapter traces this history, from empirical measures to theoretical axioms, through copulas and stress testing.

\begin{intuition}
A risk measure answers the question: ``How much can I lose at most, with a
given probability, over a given horizon?'' VaR answers with a quantile, but
says nothing about the severity of losses beyond that threshold. Expected
Shortfall fills this gap by averaging losses in the tail of the distribution.
\end{intuition}

\section{Value at Risk (VaR)}

\begin{definition}[Value at Risk]
\index{Value at Risk}
Let $L$ be the random variable representing the portfolio loss over a horizon
$\Delta t$. The \textbf{Value at Risk}\index{VaR} at confidence level
$\alpha \in (0, 1)$ is
\[
    \mathrm{VaR}_\alpha(L) = \inf\{l \in \R : \PP(L > l) \leq 1 - \alpha\}
    = F_L^{-1}(\alpha),
\]
where $F_L^{-1}$ denotes the generalised quantile of $L$.
\end{definition}

\begin{example}[Gaussian VaR]
If $L \sim \mathcal{N}(\mu, \sigma^2)$, then
$\mathrm{VaR}_\alpha = \mu + \sigma\,\Phi^{-1}(\alpha)$,
where $\Phi$ is the standard normal CDF. For $\alpha = 0.99$:
$\mathrm{VaR}_{99\%} = \mu + 2.326\,\sigma$.
\end{example}

\begin{proposition}[VaR computation methods]
\begin{enumerate}
    \item \textbf{Historical}: $\mathrm{VaR}_\alpha =$ empirical quantile of
    past returns.
    \item \textbf{Parametric (variance-covariance)}: assumes a distribution
    (often Gaussian) and uses estimated $\mu$, $\sigma$.
    \item \textbf{Monte Carlo}: simulate scenarios and compute the quantile
    of simulated losses.
\end{enumerate}
\end{proposition}

\begin{attention}[title=\textbf{Limitations of VaR}]
VaR does not satisfy sub-additivity in general: one can have
$\mathrm{VaR}_\alpha(L_1 + L_2) > \mathrm{VaR}_\alpha(L_1) + \mathrm{VaR}_\alpha(L_2)$.
This means diversification can \emph{appear} to increase risk, which is a
major conceptual flaw.
\end{attention}

\section{Expected Shortfall (CVaR)}

\begin{definition}[Expected Shortfall]
\index{Expected Shortfall}\index{CVaR}
The \textbf{Expected Shortfall} (or CVaR, Conditional Value at Risk) at
level $\alpha$ is
\[
    \mathrm{ES}_\alpha(L) = \frac{1}{1-\alpha}\int_\alpha^1
    \mathrm{VaR}_u(L)\,\mathrm{d}u.
\]
If $F_L$ is continuous,
$\mathrm{ES}_\alpha(L) = \E[L \mid L \geq \mathrm{VaR}_\alpha(L)]$.
\end{definition}

\begin{example}[Gaussian ES]
If $L \sim \mathcal{N}(\mu, \sigma^2)$:
\[
    \mathrm{ES}_\alpha = \mu + \sigma\,\frac{\phi(\Phi^{-1}(\alpha))}{1 - \alpha},
\]
where $\phi$ is the standard normal density. For $\alpha = 0.99$:
$\mathrm{ES}_{99\%} \approx \mu + 2.665\,\sigma$.
\end{example}

\begin{theorem}[Dual representation of ES]
\index{dual representation}
\[
    \mathrm{ES}_\alpha(L) = \sup_{Q \in \mathcal{Q}_\alpha} \E_Q[L],
\]
where $\mathcal{Q}_\alpha = \{Q \ll \PP : \mathrm{d}Q/\mathrm{d}\PP \leq
1/(1-\alpha)\}$. This shows that ES is the worst-case expectation over an
ambiguity set.
\end{theorem}

\section{Coherent Risk Measures}

\begin{definition}[Coherent risk measure]
\index{coherent risk measure}
A functional $\rho : \mathcal{X} \to \R$ is a \textbf{coherent risk measure}
(Artzner, Delbaen, Eber, Heath, 1999) if it satisfies:
\begin{enumerate}
    \item \textbf{Monotonicity}: $L_1 \leq L_2$ a.s. $\implies$
    $\rho(L_1) \leq \rho(L_2)$.
    \item \textbf{Translation invariance}: $\rho(L + c) = \rho(L) + c$ for
    all $c \in \R$.
    \item \textbf{Positive homogeneity}: $\rho(\lambda L) = \lambda\rho(L)$
    for all $\lambda > 0$.
    \item \textbf{Sub-additivity}: $\rho(L_1 + L_2) \leq \rho(L_1) + \rho(L_2)$.
\end{enumerate}
\end{definition}

\begin{theorem}[ADEH, 1999]
\begin{enumerate}
    \item Expected Shortfall is a \textbf{coherent} risk measure.
    \item VaR is \textbf{not} coherent (it violates sub-additivity).
    \item Every coherent risk measure admits a representation
    $\rho(L) = \sup_{Q \in \mathcal{Q}} \E_Q[L]$ for a closed convex set
    $\mathcal{Q}$ of probability measures.
\end{enumerate}
\end{theorem}

\begin{keyformulas}[title=\textbf{VaR vs ES Comparison}]
\begin{center}
\begin{tabular}{lcc}
\textbf{Property} & \textbf{VaR} & \textbf{ES} \\
\hline
Sub-additivity & No & Yes \\
Coherence & No & Yes \\
Tail sensitivity & No & Yes \\
Basel III+ regulation & No & Yes \\
Ease of backtesting & Yes & Difficult \\
\end{tabular}
\end{center}
\end{keyformulas}

\section{Stress Testing}

\begin{definition}[Stress test]
\index{stress test}
A \textbf{stress test} evaluates portfolio losses under predetermined or
historical extreme scenarios:
\[
    L_{\text{stress}} = \sum_{i=1}^n \Delta_i \cdot S_i^{\text{stress}},
\]
where $\Delta_i$ are sensitivities and $S_i^{\text{stress}}$ are shocks
applied to risk factors.
\end{definition}

\begin{proposition}[Types of stress tests]
\begin{enumerate}
    \item \textbf{Historical}: reproduce past crises (2008, 2020).
    \item \textbf{Hypothetical}: constructed scenarios (300 bp rate shock,
    sector collapse).
    \item \textbf{Reverse stress tests}: find the scenario leading to a given
    critical loss.
\end{enumerate}
\end{proposition}

\section{Copulas and Dependence Modelling}

\begin{definition}[Copula]
\index{copula}
A \textbf{copula} $C : [0,1]^d \to [0,1]$ is a joint CDF with uniform
marginals on $[0,1]$. It captures the dependence structure independently of
the marginal distributions.
\end{definition}

\begin{theorem}[Sklar, 1959]
\index{Sklar's theorem}
For any joint distribution $F$ with marginals $F_1, \ldots, F_d$, there
exists a copula $C$ such that
\[
    F(x_1, \ldots, x_d) = C(F_1(x_1), \ldots, F_d(x_d)).
\]
If the $F_i$ are continuous, $C$ is unique.
\end{theorem}

\begin{example}[Gaussian copula]
\index{copula!Gaussian}
The \textbf{Gaussian copula} with correlation matrix $\Sigma$ is
\[
    C_\Sigma^{\mathrm{Ga}}(u_1, \ldots, u_d)
    = \Phi_\Sigma(\Phi^{-1}(u_1), \ldots, \Phi^{-1}(u_d)),
\]
where $\Phi_\Sigma$ is the multivariate normal CDF.
\end{example}

\begin{definition}[Clayton copula]
\index{copula!Clayton}
The \textbf{Clayton copula} (bivariate) is
\[
    C_\theta(u, v) = (u^{-\theta} + v^{-\theta} - 1)^{-1/\theta},
    \quad \theta > 0.
\]
It models strong lower tail dependence (co-crash risk).
\end{definition}

\begin{definition}[Tail dependence]
\index{tail dependence}
The \textbf{lower tail dependence coefficient} is
\[
    \lambda_L = \lim_{u \to 0^+} \frac{C(u, u)}{u}.
\]
For the Gaussian copula, $\lambda_L = 0$ (unless $\rho = 1$). For the Clayton
copula, $\lambda_L = 2^{-1/\theta} > 0$.
\end{definition}

\begin{remark}
The Gaussian copula underestimates tail dependence. This was one of the causes
of the 2008 subprime crisis, where CDO models used Li's (2000) Gaussian copula
without accounting for extreme dependence.
\end{remark}

\begin{codeblock}[title=\textbf{VaR and ES computation via Monte Carlo}]
\begin{minted}{python}
import numpy as np
from scipy.stats import norm

np.random.seed(42)
n_sim = 100_000
mu, sigma = 0.0, 0.02
alpha = 0.99

returns = np.random.normal(mu, sigma, n_sim)
losses = -returns

var_99 = np.percentile(losses, 100 * alpha)
es_99 = losses[losses >= var_99].mean()

print(f"VaR 99%: {var_99:.4f}")
print(f"ES 99%:  {es_99:.4f}")
print(f"Analytical VaR: {-mu + sigma * norm.ppf(alpha):.4f}")
\end{minted}
\end{codeblock}

\section{Exercises}

\begin{exercise}[Portfolio VaR]
A portfolio worth EUR 10M has daily returns
$\mathcal{N}(0.05\%, (1.5\%)^2)$. Compute the 99\% VaR at 1 day and 10 days
(square root of time rule).
\end{exercise}

\begin{exercise}[Sub-additivity violation]
Construct an explicit example of two random variables $L_1, L_2$ such that
$\mathrm{VaR}_{0.95}(L_1 + L_2) > \mathrm{VaR}_{0.95}(L_1) +
\mathrm{VaR}_{0.95}(L_2)$.
\textit{Hint}: use concentrated Bernoulli losses.
\end{exercise}

\begin{exercise}[Clayton copula]
Simulate 10,000 pairs $(U, V)$ from a Clayton copula with $\theta = 2$.
Estimate the lower tail dependence coefficient and compare with the
theoretical value $\lambda_L = 2^{-1/2}$.
\end{exercise}

\begin{exercise}[Convex risk measures]
Show that every coherent risk measure is convex. Give an example of a convex
risk measure that is not coherent (positive homogeneity violated).
\end{exercise}

\begin{exercise}[Stress test]
Design a historical stress test based on the 2008 crisis for a portfolio
containing 60\% equities and 40\% bonds. Compute the loss under this scenario
and compare with the 99\% VaR.
\end{exercise}

\begin{exercise}[Backtesting]
Explain the Kupiec test for VaR backtesting. If over 250 trading days, 5 VaR
exceedances are observed at the 99\% level, is the VaR correctly calibrated
at the 5\% significance level?
\end{exercise}
